\documentclass[12pt]{article}

\input{../../_assets/preambulo.tex}


\begin{document}

    % 1. Foto de fondo
    % 2. Título
    % 3. Encabezado Izquierdo
    % 4. Color de fondo
    % 5. Coord x del titulo
    % 6. Coord y del titulo
    % 7. Fecha

    
    \input{../../_assets/portada}
    \portadaExamen{ffccA4.jpg}{Modelos\\Matemáticos I\\Examen V}{MM I. Examen V}{MidnightBlue}{-8}{28}{2026}{Roxana Acedo Parra}

    \begin{description}
        \item[Asignatura] Modelos Matemáticos I.
        \item[Curso Académico] 2025-26.
        \item[Grado] Doble Grado en Ingeniería Informática y Matemáticas.
        \item[Grupo] Único.
        \item[Profesor] María José Cáceres Granados.
        \item[Descripción] Prueba 2.
        \item[Fecha] 26 de mayo de 2026.
        \item[Duración] 105 minutos.
    
    \end{description}
    \newpage
    
    \begin{ejercicio}[4 puntos]Para cada valor $0 \ne b \in \bb{R}$, determina el comportamiento asintótico de las soluciones de
        \begin{equation}
            x_{n+2} - 2x_{n+1} + ax_n = b,
        \end{equation}
        en función del parámetro $a \in \bb{R}$.
    \end{ejercicio}

    \begin{ejercicio}Un grupo de ecología está estudiando la posibilidad de emplear el modelo de Leslie para sus estudios. Para ello está analizando el comportamiento del modelo con una matriz sencilla, para después pasar a una de dimensión mayor, y con coeficientes extraídos de sus observaciones. La matriz que está estudiando es
    $$M = \begin{pmatrix}
    0 & 0 & 1 \\
    1 & 0 & 0 \\
    0 & 1 & 0
    \end{pmatrix}$$
        De forma razonada, responde a las siguiente preguntas:
        \begin{enumerate}[label=\alph*)]
            \item \textit{(1 punto)} ¿Qué ocurrirá con la población a largo plazo, en cuanto al número de hembras en cada grupo de edad y la pirámide de edad?
            \item \textit{(1 punto)} ¿Cambia el comportamiento a largo plazo en función del dato inicial?
            \item \textit{(1 punto)} Según este modelo, ¿cómo sería el comportamiento a largo plazo de la población si las probabilidades de pasar de un grupo al siguiente fuesen 0,5?
        \end{enumerate}
    \end{ejercicio}

    \begin{ejercicio} Se pretende estudiar cómo se propaga una enfermedad en el seno de una comunidad de tamaño $N > 0$. Para ello, se considera que dicha comunidad está estructurada en los siguientes tres grupos de población:
        \begin{itemize}
            \item $I = $ personas infectadas
            \item $S = $ personas susceptibles de estarlo
            \item $R = $ personas recuperadas
        \end{itemize}
        y se propone el siguiente modelo:
        \begin{equation}
        \begin{aligned}
            S_{n+1} &= S_n - aS_nI_n \\
            I_{n+1} &= (1-b)I_n + aS_nI_n \\
            R_{n+1} &= R_n + bI_n
        \end{aligned}
        \end{equation}
        donde $a > 0$ y $b > 0$ representan el ritmo de contagio y el de recuperación, respectivamente. Asimismo, se cumple que
        
        $$S_0 + I_0 + R_0 = N$$
        
        donde $S_0 \ge 0$, $I_0 \ge 0$ y $R_0 \ge 0$ representan el número inicial de personas de cada grupo. Se pide:
        \begin{enumerate}[label=\alph*)]
            \item \textit{(1 punto)} Determina el tipo de interacción que experimentan entre sí los grupos de población $S$ e $I$.
            \item \textit{(1 punto)} Calcula los puntos de equilibrio del modelo con sentido biológico.
            \item \textit{(1 punto)} Para el caso en que $N = 10$, $a = 0{,}01$ y $b = 0{,}025$, estudia la estabilidad del punto de equilibrio ($S = 5$, $I = 0$, $R = 5$).
        \end{enumerate}
    \end{ejercicio}
\end{document}